\documentclass[12pt,a4paper]{amsart}
\usepackage[margin=1in]{geometry}
\usepackage{mathtools,amssymb}
\mathtoolsset{showonlyrefs}
\numberwithin{equation}{section}
\usepackage[hidelinks]{hyperref}
\usepackage{amsrefs}

\newtheorem{theorem}{Theorem}[section]
\newtheorem{lemma}[theorem]{Lemma}
\newtheorem{proposition}[theorem]{Proposition}
\newtheorem{corollary}[theorem]{Corollary}
\newtheorem{claim}[theorem]{Claim}
\theoremstyle{definition}
\newtheorem{definition}[theorem]{Definition}
\newtheorem{example}[theorem]{Example}
\theoremstyle{remark}
\newtheorem{remark}[theorem]{Remark}

% Replace \PaperTitle with the paper's title or define it in a local source.
% Leave the author byline empty.
\title{\PaperTitle}

\begin{document}
% If using an abstract, put the abstract environment before \maketitle.
\maketitle
% In the introduction, develop the problem and method, state all target
% results precisely, then discuss related literature.
% Write the full proofs in the main text or appendices. Use \eqref for formulas.
% For a proof ending in prose, close the proof immediately after its last sentence.
% For a final display, place \qedhere on the final equation row where appropriate.
% Check the square on the rendered last line, including numbers and line breaks.
% Text, appendices, and TeX resources can use local \input files.
% When citing sources, put verified \bib entries in bibdiv/biblist in this file.
\end{document}
